% Ad Familiares 13.26 --- headnote
% ad-familiares-13-26, 46 BC, Rome

\textit{Cicero to Servius Sulpicius Rufus, proconsul of
Achaia, written from Rome in 46 BC (the manuscript
dateline: \textit{Scr. Romae, ut videtur, a. 708 (46)}).
The longest of the five Servius-recommendations in this
sub-sequence, and the most substantively legal. The
principal is L. Mescinius Rufus, who had been Cicero's
quaestor in Cilicia in 51--50 BC --- the
\textit{quaestoria necessitudo} between proconsul and
quaestor being one of the formal Roman bonds that
followed both men for life. Mescinius has come into an
inheritance, in Achaia, from his late brother M. Mindius,
a Roman trader at Elis; certain heirs or creditors are
likely to resist, and Cicero is asking Servius to use
both his judicial power and his personal weight to
unravel the business.}

\textit{The letter is a small masterpiece of
provincial-legal recommendation, and the
\textit{non-vulgaris} register-shift is visible. Three
moves are worth noting. First, Cicero pre-loads the
request by attaching a covering letter from M. Lepidus
(the consul of 46), explicitly described as not a
directive --- "not such as enjoin anything upon you (for
we do not consider that to be consistent with your
standing), but, so to say, in the manner of letters of
recommendation" --- a finely calibrated face-saving
formula for asking the governor to do something while
denying that he is being asked. Second, Cicero offers
Servius a Roman-forum escape: any party who proves
recalcitrant may be remitted to Rome on the ground that
the case involves a senator. Third, the close fuses the
two motives Cicero typically separates --- the welfare
of the friend and the visible standing of the
recommender --- by stating both with equal weight. The
piece is professional commendation business done at its
high tier.}
